\section{Difficulties Encountered}
\label{sec:difficulties}
\textbf{Data access.} The proposal planned to collect labelled listings from resale sites, by hand or
through APIs. Automated collection turned out not to be possible: eBay's Buy APIs are now limited to
approved partners in production and its June 2025 API license update restricts using eBay content
with AI, while Grailed and Depop block automated requests. A search of Hugging Face, Kaggle, Zenodo,
academic releases and museum APIs found no public dataset of designer \emph{label} photos with era
labels (The Met's API returns no images for these designers; the V\&A has only dozens per brand),
hence the three proxies.
\textbf{What the proxies measure.} Fashion-Gen's ``era'' is the season in which SSENSE sold an item
(2014--2018), not its production date; its prices are retail, not resale; some brands occur only in
some years; and catalog or runway photos look nothing like a label close-up. Low-level photographic
cues also carry era information (an 8$\times$8 thumbnail predicts the Fashion-Gen year 41.4\% of the
time vs.\ a 31.7\% majority), and shoot-level cues such as models and styling may too, so we treat
proxy era accuracies as optimistic.
\textbf{Data quality and leakage.} SSENSE reuses one outfit photo for every product in a look, and a
review of our own code found that our first, product-grouped split therefore put 60 test images whose
exact copies were in training (inflating $k$-NN most). We now merge products that share an identical
or near-identical photo, rebuilt the splits and re-ran every Fashion-Gen experiment; a leakage audit
finds no validation or test image with an exact, near-duplicate or CLIP-cosine $\geq$0.99 copy in
training on any of the three datasets. In the second-hand data, 46\% of garments have no usable brand
annotation, some ``label'' photos show care tags, and many labels are upside down or sideways.
\textbf{Licensing and compute.} Fashion-Gen is research-only and Europeana catwalk images are
\emph{In Copyright}; only the second-hand data is openly licensed (CC BY 4.0), so we share code but
never images. Everything runs on one Apple-silicon laptop; fine-tuning CLIP ViT-B/16 takes about
9 minutes per epoch, so runs are limited to 8 epochs and a single seed for now.

\section{Changes to the Proposal}
\label{sec:changes}
\textbf{Data.} Development now happens on public proxy datasets. The target data come from listing
photos we save by hand (label close-up, title, era claim and sale price) and from our own in-store
phone photos, which remain the held-out real-world test set.
\textbf{Era labels.} Eras are defined per brand by visible label changes and printed date codes
(Sec.~\ref{sec:eras}) instead of fixed seasons; neighbouring eras are merged where we cannot collect
at least 50 label photos.
\textbf{Value.} Retrieval-based price estimation is implemented and evaluated on retail prices; it
will be re-run on resale prices from the listings we collect. Automated ``web search for current
listings'' is replaced by these curated comparables, because the marketplace APIs are closed to us.
\textbf{Brand scope.} The pipeline handles any set of houses. The final set (CDG, Margiela,
Undercover, plus Yohji Yamamoto or Junya Watanabe if data allow) depends on how many label photos we
can collect per era.
\textbf{Updated timeline.} Weeks 6--7: collect in-store and listing label photos ($\geq$50 per
brand-era class); extend the OCR module to printed date codes; train frozen and fine-tuned models on
the real data. Week 8: evaluate OCR--image fusion for brand and date codes on the real label photos, value
estimation on resale comparables, error analysis. Week 9:
final evaluation on the held-out in-store set, webpage, presentation, report.
